    \documentclass[11pt,a4paper]{article}
    \usepackage[utf8]{inputenc}
    \usepackage[T1]{fontenc}
    \usepackage{amsmath, amssymb, amsthm}
    \usepackage{geometry}
    \usepackage{booktabs}
    \usepackage{array}
    \usepackage{parskip}
    \usepackage{enumitem}

    \geometry{margin=2.5cm}
    \newcolumntype{L}[1]{>{\raggedright\arraybackslash}p{#1}}

    \begin{document}

    \section{Starting Formula: Multi-Asset Objective}

    \subsection*{Discrete form (Eq.~15 built from Eqs.~29 and 30)}

    \begin{equation}
    U[\mathbf{x}]
    = \sum_{j=1}^{M}\left[
        \tfrac{1}{2}\gamma_j X_j^2
        + \varepsilon_j \sum_{k=1}^{N}|n_{jk}|
        + \frac{\tilde{\eta}_j}{\tau}\sum_{k=1}^{N} n_{jk}^2
        \right]
    + \lambda \sum_{k=1}^{N} \tau\,\mathbf{x}_k^\top C\,\mathbf{x}_k
    \end{equation}

    \textbf{When used.} The total quantity the trader minimises: expected execution
    cost (Eq.~29) plus $\lambda$ times the variance of that cost (Eq.~30). Every later
    equation in the time-dependent formulation is derived from this one.

    \textbf{Source.} R.~Almgren and N.~Chriss, ``Optimal Execution of Portfolio
    Transactions,'' December 2000 (page numbers are those of the running headers).
    \begin{itemize}[leftmargin=1.5em,nosep]
    \item Objective $E+\lambda V$ (Eq.~15): Sec.~2.1, ``The definition of the frontier,'' p.~13.
    \item Expected cost with linear impact (Eq.~8) and $\tilde\eta=\eta-\tfrac12\gamma\tau$: Sec.~1.5, p.~11.
    \item Variance (Eq.~5): Sec.~1.4, p.~10.
    \item Multi-asset expected cost and variance (Eqs.~29--30), covariance $C=\sigma\sigma^\top$:
            Appendix~A.1, ``Trading model,'' pp.~36--37.
    \item Multi-asset optimality condition (Eq.~31), used as the final check: Appendix~A.2, p.~37.
    \end{itemize}
    \textbf{Caveats.} (1)~The paper prints Eq.~(29) for a general permanent-impact
    matrix $\Gamma$ and temporary-impact matrix $H$ (with a cross term). The
    asset-by-asset sum above is the special case of diagonal $\Gamma$ and $H$ (see
    Eq.~32 of the reference). If $\Gamma$ has an antisymmetric part, Eq.~(31) in the
    paper contains an extra term that this formulation does not have.
    (2)~The combined multi-asset objective $E[\mathbf{x}]+\lambda V[\mathbf{x}]$ is
    not printed as a single equation; it is Eq.~(15) applied to Eqs.~(29)--(30).
    (3)~Page numbers were read from a machine extraction of the PDF; verify them
    against the PDF before citing.

    \begin{center}
    \begin{tabular}{@{}L{3.8cm}L{10cm}@{}}
    \toprule
    \textbf{Term} & \textbf{Meaning} \\
    \midrule
    $U[\mathbf{x}]$ & Total objective for the whole trading schedule; the quantity to minimise \\
    $j = 1,\dots,M$ & Asset index ($M=2$ in this project) \\
    $k = 1,\dots,N$ & Time-interval index; $\tau = T/N$ \\
    $X_j$ & Total shares of asset $j$ to liquidate \\
    $n_{jk}$ & Shares of asset $j$ sold in interval $k$ \\
    $\mathbf{x}_k$ & Vector of shares still held in each asset at interval $k$ (one entry per asset) \\
    $\tfrac{1}{2}\gamma_j X_j^2$ & Permanent-impact cost of asset $j$; identical for every schedule, so it drops out of the optimisation \\
    $\varepsilon_j\sum_k |n_{jk}|$ & Bid-ask spread cost; equals $\varepsilon_j X_j$ for any monotone sell, so it also drops out \\
    $\dfrac{\tilde\eta_j}{\tau}\sum_k n_{jk}^2$ & Temporary-impact (slippage) cost; grows with the square of the trade size, so it penalises selling fast \\
    $\tilde\eta_j$ & $\eta_j - \tfrac12\gamma_j\tau$, the effective temporary-impact coefficient \\
    $\lambda$ & Risk aversion; larger $\lambda$ makes holding inventory more costly \\
    $C$ & $M\times M$ covariance matrix of asset returns per unit time \\
    $\tau\,\mathbf{x}_k^\top C\,\mathbf{x}_k$ & Variance contribution of interval $k$; holding more of correlated assets raises it \\
    $\sum_k \tau(\cdots)$ & Accumulates the variance over the whole horizon (Eq.~30) \\
    \bottomrule
    \end{tabular}
    \end{center}

    \subsection*{Continuous-time form}

    Under sales-only trading and exact liquidation, the two discrete terms that
    do not depend on the schedule have the fixed value
    $K(\mathbf X)=\sum_j(\tfrac12\gamma_jX_j^2+\varepsilon_jX_j)$. Let
    $\tau\to dt$, set $n_{jk}=\tau v_j$, and write
    $H=\operatorname{diag}(\eta_1,\dots,\eta_M)
    =\lim_{\tau\to0}\operatorname{diag}(\tilde\eta_1,\dots,\tilde\eta_M)$:

    \begin{equation}
    \begin{gathered}
    J_{\mathrm{full}}[\mathbf{v}]=K(\mathbf X)+J[\mathbf v],\\
    J[\mathbf{v}]
    = \int_0^T \Big( \mathbf{v}(t)^\top H\,\mathbf{v}(t)
        + \lambda\,\mathbf{x}(t)^\top C\,\mathbf{x}(t) \Big)\,dt,
    \\\dot{\mathbf{x}}(t) = -\mathbf{v}(t),
    \quad \mathbf{x}(0)=\mathbf{X},\quad \mathbf{x}(T)=\mathbf{0}
    \end{gathered}
    \end{equation}

    \textbf{When used.} $J_{\mathrm{full}}$ is the total cost; $J$ is its
    schedule-dependent part and is the control problem Bellman's principle turns
    into the time-dependent Hamilton--Jacobi equation. The spread offset is
    constant only when $\mathbf{x}(T)=\mathbf 0$; it need not be constant under a
    finite terminal penalty that permits residual inventory.

    \textbf{Source.} This continuous-time form is derived in these notes; it does not
    appear in Almgren and Chriss, who work in discrete time and solve by calculus of
    variations. It follows from the discrete objective above by letting $\tau\to0$.

    \begin{center}
    \begin{tabular}{@{}L{3.8cm}L{10cm}@{}}
    \toprule
    \textbf{Term} & \textbf{Meaning} \\
    \midrule
    $J_{\mathrm{full}}[\mathbf{v}]$ & Total cost, including the schedule-independent offset $K(\mathbf X)$ \\
    $J[\mathbf{v}]$ & Schedule-dependent integral minimized by the HJB control \\
    $\mathbf{v}(t)$ & Selling rate (shares per unit time) in each asset; this is the control \\
    $\mathbf{x}(t)$ & Inventory vector (the state) at time $t$ \\
    $H$ & Continuum impact matrix $\operatorname{diag}(\eta_j)$ \\
    $\mathbf{v}^\top H\,\mathbf{v}$ & Running slippage cost; large when selling fast \\
    $\lambda\,\mathbf{x}^\top C\,\mathbf{x}$ & Running risk cost; large when holding a big or concentrated position \\
    $\dot{\mathbf{x}}=-\mathbf{v}$ & Selling reduces inventory \\
    $\mathbf{x}(0)=\mathbf{X},\ \mathbf{x}(T)=\mathbf{0}$ & Start with the full position; be completely sold by the horizon $T$ \\
    \bottomrule
    \end{tabular}
    \end{center}

    \end{document}